%% kapa: exam=1 round=32 year=2021 period=1 form=A subject=economics subject_ko=경제학원론 seq=34 number=74
%% source: Q-Net 감정평가사 기출문제 (cst00309) artlSeq=5208397 file=extracted/1008/5208397/A형/2021년도 제32회 감정평가사 1차 1교시 A형.pdf
%% conversion: kordoc 4.24.51 -> tools/gen1.py (auto); review=proofread-sonnet; flags=none
\begin{question}{74}
변동환율제를 채택한 A국이 긴축재정을 실시하였다. 먼델-플레밍 모형을 이용한 정책 효과에 관한 설명으로 옳은 것을 모두 고른 것은? (단, 완전한 자본이동, 소국개방경제, 국가별 물가수준 고정을 가정한다.)
\begin{framedtext}
\makebox[0.5\linewidth][l]{ㄱ. 원화가치는 하락한다.}ㄴ. 투자지출을 증가시킨다. \par
\makebox[0.5\linewidth][l]{ㄷ. 소득수준은 변하지 않는다.}ㄹ. 순수출이 감소한다. \par
\end{framedtext}
\begin{choices}
  \item ㄱ, ㄴ
  \item ㄱ, ㄷ
  \item ㄱ, ㄹ
  \item ㄴ, ㄷ
  \item ㄴ, ㄹ
\end{choices}
\end{question}
\answer{2}
